\documentclass[10pt,a4paper]{article}
\usepackage{amsmath,amssymb,amsthm,graphicx}
\usepackage{xfrac}
% Minimal editorial layout adapter for the pinned JEP source
% sources/jep-2022-195-contact-systolic/pinned-original-native.tex.
% The publisher document class (cedram) is not part of the retained material, and it is
% the class that used to define a handful of plain-text abbreviations and sectioning
% helpers used by the author's prose. This adapter supplies exactly those and nothing
% else: no mathematics, no notation, no theorem environment, no counter.
% CC BY 4.0, attribution to the original authors.
\providecommand{\ie}{i.e.\ }
\providecommand{\eg}{e.g.\ }
\providecommand{\skpt}{\smallskip}
\providecommand{\Subsection}{\subsection}
\providecommand{\Subsubsection}{\subsubsection}
\providecommand{\loccit}{\emph{loc.\ cit.}}
\providecommand{\dpl}{\displaystyle}
\providecommand{\backmatter}{}
\usepackage[scr=rsfs,cal=euler]{mathalfa}
\let\mathds\mathbb
\let\epsilon\varepsilon
\newcommand{\Changel}{\mathcode`l="7160}
\newcommand{\Changelback}{\mathcode`l="716C}
\newcommand{\NoChangel}[1]{%
\expandafter\let\csname old\string#1\endcsname=#1
\let#1=\relax
\newcommand{#1}{\mathcode`l="716C\csname old\string#1\endcsname\mathcode`l="7160 }%
}

\NoChangel{\log}
\NoChangel{\ln}
\NoChangel{\lim}
\NoChangel{\limsup}
\NoChangel{\liminf}
\NoChangel{\varlimsup}
\NoChangel{\varliminf}
\NoChangel{\varinjlim}
\NoChangel{\varprojlim}

\newenvironment{enumeratei}
{\bgroup\def\theenumi{\roman{enumi}}\def\theenumii{\arabic{enumii}}\begin{enumerate}}
{\end{enumerate}\egroup}
\newcommand{\Psfrac}[2]{(\sfrac{#1}{#2})}
\newcommand{\psfrac}[2]{\sfrac{(#1)}{#2}}
\newcommand{\spfrac}[2]{\sfrac{#1}{(#2)}}
\newcommand\mto{\mathchoice{\longmapsto}{\mapsto}{\mapsto}{\mapsto}}
\newcommand\hto{\mathrel{\lhook\joinrel\to}}

\usepackage{graphicx}

\newtheorem{thm}{Theorem}[section]
\newtheorem{cor}[thm]{Corollary}
\newtheorem{lem}[thm]{Lemma}
\newtheorem{prop}[thm]{Proposition}
\newtheorem{quest}{Question}[section]
\newtheorem{conj}{Conjecture}[section]
\newtheorem{maintheorem}{Theorem}
\theoremstyle{definition}
\newtheorem{rem}[thm]{Remark}
\newtheorem{exm}[thm]{Example}
\newtheorem{ass}[thm]{Assumption}
\renewcommand*{\themaintheorem}{\Alph{maintheorem}}

\DeclareMathOperator{\area}{area}
\DeclareMathOperator{\lcm}{lcm}
\DeclareMathOperator{\const}{const}
\DeclareMathOperator{\Cal}{Cal}
\DeclareMathOperator{\Flux}{\mathrm{Flux}}
\newcommand{\id}{\mathrm{id}}
\newcommand{\Ham}{\mathrm{Ham}}
\newcommand{\std}{\mathrm{std}}
\newcommand{\vol}{\mathrm{vol}}
\newcommand{\fix}{\mathrm{fix}}
\newcommand{\pr}{\mathrm{pr}}
\newcommand{\Tan}{\mathrm{T}}
\newcommand{\N}{\mathds{N}}
\newcommand{\Z}{\mathds{Z}}
\newcommand{\R}{\mathds{R}}
\newcommand{\CC}{\mathds{C}}\let\C\CC
\newcommand{\UU}{\mathcal{U}}
\newcommand{\VV}{\mathcal{V}}

\let\eul e
\let\diff d
\title{Every exact symplectomorphism of a compact surface with prescribed collar behaviour that is C1-close to the identity is generated by a quasi-autonomous Hamiltonian, with C1-small Hamiltonian and a collar expansion b\_C + r\textasciicircum{}2 h\_C}
\author{Alberto Abbondandolo, Christian Lange, Marco Mazzucchelli}
\date{doi:10.5802/jep.195; CC BY 4.0}
\begin{document}
\maketitle
\setcounter{section}{2}
\setcounter{subsection}{2}
\setcounter{thm}{6}
\noindent\textit{Editorial note.} The counted unit is the author's labelled theorem t:quasi\_autonomous at source lines 584-598 with the proof that occupies source section 2.3, namely the retained block at source lines 693-810 and the author's proof environment at source lines 811-959; the author states at source line 693 that this section is the proof of the theorem, and the Weinstein tubular neighbourhood lemma proved inside it is invoked by its label inside the retained proof.\par\medskip
\noindent\textit{Editorial note.} The retained counter value for the theorem-class counter makes the counted theorem print as 2.8, which is its number in the pinned source; the retained standing Assumption then prints as 2.7 instead of 2.4, because the numbered items of source section 2 between the Assumption and the theorem are not retained. The reference to the Assumption inside the counted statement resolves to the retained block, so the excerpt is internally consistent.\par\medskip
\noindent\textit{Editorial note.} The retained preamble is the author's own macro, operator, environment and theorem block of the pinned source; the publisher document class cedram and its submission-date commands are excluded, and no mathematical macro, symbol, statement or proof step was added, changed or removed.\par\medskip
\noindent\textit{Editorial note.} The context spans are the author's standing Assumption and the remark on applications of the theorem; they are marked as context and are not counted as this task's proof.\par\medskip
\noindent\textit{Editorial note.} The retained bibliography is the article's own reference list, flattened at the end of the document so that the citations of the retained proof resolve; no entry was rewritten.\par\medskip
\noindent\textit{Allowed standard prerequisites (editorial).}\begin{itemize}
\item Hamiltonian isotopies, the Calabi invariant and the group of Hamiltonian diffeomorphisms with zero flux. The statement is phrased in the group of exact symplectomorphisms isotopic to the identity. The preliminaries are not re-proved; the quasi-autonomous generating Hamiltonian and the collar expansion of the theorem are the retained obligations.
\item The Weinstein tubular neighbourhood theorem on the diagonal of a surface. The counted proof invokes this lemma by its label. Its complete author proof is part of the retained material, so nothing is borrowed unproved.
\item Moser's trick and the implicit function theorem on compact surfaces with boundary. Only the standard C1-openness of the symplectic condition and the standard local inversion are assumed; every estimate that produces the C1-small quasi-autonomous Hamiltonian is retained.
\end{itemize}\medskip
\begin{ass}
\label{assumption}
Every connected component $C$ of the boundary $\partial \Sigma$ has a collar neighborhood $A_C\subset \Sigma$ and an identification $A_C\equiv [0,\rho) \times S^1$, for some $\rho>0$ such that
\[
\omega|_{A_C} = -r\, dr\wedge ds.
\]
Here, we are identifying $S^1$ with $\R/\Z$, and $(r,s)$ denotes a point in $[0,\rho) \times S^1$. Note that the orientation of $\partial \Sigma$ as boundary of the oriented surface $(\Sigma,\omega)$ coincides, under the above identification of each component $C\subset \partial \Sigma$ with $\{0\} \times S^1$, with the orientation given by $ds$.
\hfill\qed
\end{ass}

\begin{thm}
\label{t:quasi_autonomous}
Assume that the exact two-form $\omega$ on the compact surface $\Sigma$ is symplectic on $\mathrm{int}(\Sigma)$ and satisfies Assumption \ref{assumption}. Let $\phi:\Sigma\to\Sigma$ be an exact symplectomorphism that is sufficiently $C^1$-close to the identity. Then there exists a Hamiltonian isotopy $\{\phi_t\}$ from $\mathrm{id}$ to $\phi$ whose generating Hamiltonian $H_t$ is quasi-autonomous.
Moreover, for every $\epsilon>0$ there exists $\delta>0$ such that, if $\mathrm{dist}_{C^1}(\phi,\mathrm{id})<\delta$, then:
\begin{enumeratei}
\item\label{t:quasi_autonomousi}
$\|H_t\|_{C^1}<\epsilon$ and $\mathrm{dist}_{C^1}(\phi_t,\mathrm{id})<\epsilon$ for every $t\in [0,1]$;
\item\label{t:quasi_autonomousii}
in the collar neighborhood $A_C\equiv [0,\rho)\times S^1$ of each boundary component $C$ of~$\Sigma$ as in Assumption \ref{assumption}, the Hamiltonian $H_t$ has the form
\[
H_t(r,s) = b_C + r^2 h_C(t,r,s)
\]
for some real number $b_C$ and some smooth function $h_C: [0,1] \times A \to \R$ such that $|b_C|<\epsilon$ and $\|h_C\|_{C^0} < \epsilon$.
\end{enumeratei}
\end{thm}

\begin{rem}
Theorem \ref{t:quasi_autonomous} has other interesting applications. For instance, it implies that the identity in $\Ham(\Sigma,\omega)$ has a $C^1$-neighborhood on which the Hofer metric is flat. See \cite{Bialy:1994} for more about this in the setting of compactly supported Hamiltonian diffeomorphisms of $\R^{2n}$ and \cite{Lalonde:1995} for the case of compactly supported Hamiltonian diffeomorphisms of more general symplectic manifolds.
\hfill\qed
\end{rem}

The aim of this section is to prove Theorem \ref{t:quasi_autonomous}. The proof closely follows the argument of \cite[\S 4]{Benedetti:2021aa}, but for sake of completeness we work out the details.

By Assumption \ref{assumption} and up to reducing the positive number $\rho$ appearing there, we~can find a primitive $\nu_0$ of $\omega$ which on the collar neighborhood $A_C\equiv [0,\rho)\times S^1$ of each component $C$ of the boundary of $\Sigma$ has the form
\begin{equation}
\label{e:lambda|_A}
\nu_0|_{A_C} = \bigl( a_C - \tfrac{1}{2} r^2 \bigr)\, ds,
\end{equation}
for some $a_C\in \R$. Indeed, if $\nu$ is any primitive of $\omega$ and $a_C$ is its integral on the boundary component $C = \{0\} \times S^1$, then the one-form above differs from $\nu|_{A_C}$ by the differential of a function $g_C$. By adding to $\nu$ the differential of a function on $\Sigma$ that agrees with $g_C$ on a slightly reduced collar neighborhood of every component $C$ of $\partial \Sigma$, we obtain the desired primitive $\nu_0$.

For $i=1,2$, we consider the one-forms $\nu_i:=\pr_i^*\nu_0$ on $\Sigma\times\Sigma$, where
\[
\pr_i:\Sigma\times\Sigma\to\Sigma, \qquad \pr_i(z_1,z_2)=z_i,
\]
and the standard Liouville form $\lambda_{\mathrm{std}}$ on the cotangent bundle $\Tan^*\Sigma$, which is uniquely defined by the equation $\alpha^*\lambda_{\mathrm{std}}=\alpha$ for all one-forms $\alpha$ on $\Sigma$.

If $z_1$ and $z_2$ are points on the same connected component $C$ of $\partial \Sigma$, then they are identified with pairs $(0,s_1)$, $(0,s_2)$ in $A_C=[0,\rho)\times S^1$. When $s_2$ and $s_1$ are not antipodal in $S^1=\R/\Z$, meaning that $|s_2-s_1|< \sfrac 12$ for suitable lifts to $\R$, we denote by $[z_1,z_2]$ the unique shortest oriented arc in $C$ from $z_1$ to $z_2$, so that
\[
\int_{[z_1,z_2]} \!\!\! \diff s < \frac12.
\]

The next result is a version of Weinstein tubular neighborhood theorem in our setting.

\begin{lem}[Weinstein tubular neighborhood]
\label{l:Weinstein}
There exists an open neighborhood $U\subset\Sigma\times\Sigma$ of the diagonal $\Delta_\Sigma=\{(z,z)\mid z\in\Sigma\}$, an open neighborhood \hbox{$V\subset\Tan^*\Sigma$} of the $0$-section $0_\Sigma\subset\Tan^*\Sigma$, and a smooth map $\psi:U\to V$ that restricts to a diffeomorphism $\psi:U\cap\mathrm{int}(\Sigma\times\Sigma)\to V\cap \Tan^* \mathrm{int} (\Sigma)$, and satisfies
\[
\psi(\Delta_\Sigma)=0_\Sigma,
\qquad
\psi^*\lambda_{\std}=\nu_2-\nu_1 + \diff f,
\] where $f:\Sigma\times\Sigma\to\R$ is a smooth function such that $f|_{\Delta_\Sigma}\equiv0$ and
\[
\diff f(z_1,z_2)=(\nu_1-\nu_2)_{(z_1,z_2)},
\qquad
f(z_1,z_2)=-\int_{[z_1,z_2]} \!\!\!\!\!\!\nu_0,
\]
for every pair $(z_1,z_2)\in U\cap(\partial\Sigma\times\partial\Sigma)$.
If $A_C\equiv[0,\rho)\times S^1$ is the collar neighborhood of a connected component $C$ of $\partial\Sigma$ on which $\nu_0$ has the form \eqref{e:lambda|_A}, the restriction $\psi|_{U\cap(A_C\times A_C)}$ has the form
\begin{equation}
\label{psi}
\begin{split}
\psi(r,s,R,S)=\big( R,s,\,R\,(S-s),\tfrac 1 2 (r^2-R^2) \big),
\\
\forall (r,s,R,S)\in U\cap(A_C\times A_C).
\end{split}
\end{equation}
\end{lem}

\begin{proof}
We first provide the construction within the collar neighborhood $A=A_C$ of each connected component $C$ of $\partial\Sigma$.
We consider a small enough neighborhood $W\subset A\times A$ of the diagonal $\Delta_A=\{(z,z)\mid z\in A\}$ so that, for each $(r,s,R,S)\in W$, the points $s,S\in S^1$ are not antipodal. We define the map
\[
\kappa_0:W\to \underbrace{\Tan^*A}_{A\times \R^2},
\qquad
\kappa_0(r,s,R,S)=\big( R,s,\,R\,(S-s),\tfrac 12 (r^2-R^2) \big).
\]
This map restricts to a diffeomorphism onto its image
\[
\kappa_0:W\cap\mathrm{int}(A\times A)\to \mathrm{int}(\Tan^*A),
\]
and satisfies $\kappa_0(\Delta_A)=0_A$ and
\begin{align*}
\kappa_0^*\lambda_{\std}
&=
R\,(S-s)\,\diff R + \tfrac 12(r^2-R^2)\,\diff s\\
&=
-\tfrac 12\,R^2\,\diff S + \tfrac 12\,r^2\,\diff s + \diff\big( \tfrac 12(S-s)\,R^2 \big)\\
&=
-\nu_1+\nu_2+\diff\big( (a_C -\tfrac 12 R^2)(s-S) \big)\\
&=-\nu_1+\nu_2+\diff f_0,
\end{align*}
where
\[
f_0:W\to\R,\qquad f_0(r,s,R,S)=(a_C -\tfrac 12 R^2)(s-S).
\]
Notice that $f_0|_{\Delta_A}\equiv0$ and
\begin{align*}
f_0(0,s,0,S) & =a_C (s-S) = - \int_{[z_1,z_2]} \nu_0,
\\
\diff f_0(0,s,0,S) & = a_C \,\diff s-a_C\,\diff S = \nu_1-\nu_2,
\end{align*}
where $z_1=(0,s)$ and $z_2=(0,S)$.

For some sufficiently small neighborhood $U\subset\Sigma\times\Sigma$ of the diagonal $\Delta_\Sigma$, we choose an arbitrary smooth function $f_1:U\to\R$ that coincides with $f_0$ on a neighborhood of $U\cap(\partial\Sigma\times\partial\Sigma)$ and satisfies $f_1|_{\Delta_\Sigma}\equiv0$ and $\diff f_1=\nu_1-\nu_2$ at all points of the diagonal $\Delta_\Sigma$. Up to further shrinking the neighborhood $U$, we also choose a smooth map $\kappa_1:U\to\Tan^*\Sigma$ that coincides with $\kappa_0$ on a neighborhood of $W\cap(\partial\Sigma\times\partial\Sigma)$, restricts to a diffeomorphism onto its image $\kappa_1:U\cap\mathrm{int}(\Sigma\times\Sigma)\to\Tan^*(\mathrm{int}(\Sigma))$, and such that $\kappa_1(\Delta_\Sigma)=0_\Sigma$.

We now conclude the proof by means of a typical application of Moser's trick. We~set
\[
\mu_t:=t\kappa_1^*(\lambda_{\std}) + (1-t)(\nu_2-\nu_1),
\]
and we look for an isotopy $\phi_t:U\to\Sigma\times\Sigma$, defined after possibly further shrinking the neighborhood $U$, such that $\phi_0=\id$ and $\phi_t^*\mu_t-\mu_0$ is exact. We denote by $X_t$ the time-dependent vector field generating $\phi_t$ and compute
\begin{align*}
\tfrac{\diff}{\diff t}\phi_t^*\mu_t
&=\phi_t^* ( L_{X_t} \mu_t + \tfrac{\diff}{\diff t} \mu_t) = \phi_t^*\big( X_t\,\lrcorner\,\diff\mu_t + \diff(\mu_t(X_t))+ \mu_1-\mu_0 \big)\\
&=\phi_t^*\big( X_t\,\lrcorner\,\diff\mu_t + \kappa_1^*\lambda_{\std}-\nu_2+\nu_1-\diff f_1 + \diff(\mu_t(X_t)+f_1)\big).
\end{align*}
Notice that the symplectic forms $\kappa_1^*\diff\lambda_{\std}$ and $-\diff\nu_1+\diff\nu_2$ define the same orientation, since they coincide on a neighborhood of $U\cap(\partial\Sigma\times\partial\Sigma)$. Therefore $\diff\mu_t$ is symplectic away from $\partial(\Sigma\times \Sigma)$ for every $t\in[0,1]$. We choose the vector field $X_t$ so that
\[
X_t\,\lrcorner\,\diff\mu_t + \kappa_1^*\lambda_{\std}-\nu_2+\nu_1-\diff f_1
=
0.
\]
Notice that $X_t$ vanishes on the diagonal $\Delta_\Sigma$ and on a neighborhood of $U\cap(\partial\Sigma\times\partial\Sigma)$.
Moreover
\[
\tfrac{\diff}{\diff t}\phi_t^*\mu_t
=
\phi_t^*\diff(\mu_t(X_t)+f_1).
\]
Up to shrinking the neighborhood $U$, we obtain a well defined isotopy $\phi_t:U\to\Tan^*\Sigma$ that coincides with $\kappa_0$ on a neighborhood of $U\cap(\partial\Sigma\times\partial\Sigma)$ and satisfies
\[
\phi_1^*\kappa_1^*(\lambda_{\std})=-\nu_1+\nu_2 + \diff f,
\]
where
\[
f(w)=\int_0^1 \big(\mu_t(X_t)+f_1\big)\circ\phi_t(w)\,\diff t.
\]
The desired map is $\psi:=\kappa_1\circ\phi_1$.
\end{proof}


\begin{proof}[Proof of Theorem~\ref{t:quasi_autonomous}]
We still work with the special primitive $\nu_0$ of $\omega$ satisfying \eqref{e:lambda|_A}.
By assumption, the diffeomorphism $\phi: \Sigma \to \Sigma$ satisfies
\[
\phi^* \nu_0 - \nu_0 = da
\]
for some smooth function $a$ on $\Sigma$. Note that $a$ is $C^1$-small when $\phi$ is $C^1$-close to the identity. We consider the associated map
\[
\Phi: \Sigma \to \Sigma \times \Sigma, \qquad \Phi(z) = (z,\phi(z)).
\]
We require $\phi$ to be sufficiently $C^1$-close to the identity so that the image of $\Phi$ is contained in the domain of the map $\psi:U\to\Tan^*\Sigma$ provided by Lemma~\ref{l:Weinstein}, and the image of $\psi\circ\Phi$ is a section of the cotangent bundle $\Tan^*\Sigma$. Namely, if we denote by $\pi:\Tan^*\Sigma\to\Sigma$ the projection onto the base of the cotangent bundle, the map
\[
\widetilde\phi:\Sigma\to\Sigma,
\qquad
\widetilde\phi(z)=\pi\circ\psi\circ\Phi(z)
\]
is a diffeomorphism.
We consider the smooth function $f:\Sigma\times\Sigma\to\R$ provided by Lemma~\ref{l:Weinstein}.
Since
\[
(\psi\circ\Phi)^*\lambda_{\std}
=
\Phi^*(\nu_2-\nu_1+\diff f)
=
\phi^*\nu-\nu+\diff(f\circ\Phi)
=
\diff( a+f\circ\Phi),
\]
we have that
\begin{equation}
\label{defF}
\psi\circ\Phi(z)=(\widetilde\phi(z),\diff F(\widetilde\phi(z))),
\end{equation}
where $F:\Sigma\to\R$ is the smooth generating function
\[
F(w)=( a+f\circ\Phi)\circ\widetilde\phi^{-1}(w).
\]
Identity \eqref{defF} implies that $F$ is $C^2$-small when $\phi$ is $C^1$-close to the identity.

Consider now the collar neighborhood $A=A_C\equiv[0,\rho)\times S^1$ of a connected component $C$ of $\partial\Sigma$ as in \eqref{e:lambda|_A}. For all $(r,s)\in A\cap\phi^{-1}(A)$, if we set $(R,S)=\phi(r,s)$, we~have $\widetilde\phi(r,s)=(R,s)$ and
\begin{equation}
\label{genfun}
R\,(S-s) = \partial_R F(R,s),\qquad
\tfrac 12 (r^2-R^2) = \partial_s F(R,s).
\end{equation}
This implies that $\diff F=0$ at all points of $\partial\Sigma$. In particular, $F$ is constant on each component $C$ of $\partial \Sigma$ and in a neighborhood of this component we can write $F$ as
\begin{equation}
\label{Fnearbdry}
F(R,s)=b+ R^2G(R,s),
\end{equation}
where $b=b_C$ is a real number and $G=G_C$ is a smooth function. More precisely,
\[
b= \lim_{R\to 0} \frac{\partial_R F(R,s)}{R} = S(0,s) - s
\]
and the Taylor theorem with integral remainder gives us the formula
\[
F(R,s) = b + R^2 \int_0^1 \partial_R^2 F(tR,s) (1-t)\, dt.
\]
By differentiating the first identity in \eqref{genfun} with respect to $R$, we find
\[
G(R,s) = \int_0^1 \partial_R^2 F(tR,s) (1-t)\, dt = \int_0^1 \bigl( -s + S(tR,s) + tR \,\partial_R S(tR,s) \bigr) (1-t)\, dt.
\]
The above formulas for $b$ and $G$ imply that $|b|$ and $\|G\|_{C^0}$ are both small if $\phi$ is $C^1$-close to the identity.

We now consider the isotopy
\[
\phi_t:\mathrm{int}(\Sigma)\to\mathrm{int}(\Sigma),\ t\in[0,1],
\]
with the associated map $\Phi_t:\mathrm{int}(\Sigma)\to\mathrm{int}(\Sigma\times\Sigma)$, $\Phi_t(z)=(z,\phi_t(z))$, whose image $\psi\circ\Phi_t(\mathrm{int}(\Sigma))$ is the graph of $t\,\diff F$. Notice that $\phi_t$ defines an associated diffeomorphism
\[
\widetilde\phi_t:=\pi\circ\psi\circ\Phi_t:\mathrm{int}(\Sigma)\to\mathrm{int}(\Sigma),
\]
and
\begin{equation}
\label{e:graph_tdF}
\psi\circ\Phi_t(z)=(\widetilde\phi_t(z),t\,\diff F(\widetilde\phi_t(z))).
\end{equation}
The endpoints of the isotopy are $\phi_0=\id$ and $\phi_1=\phi$. We claim that $\phi_t$ extends as a smooth isotopy $\phi_t:\Sigma\to\Sigma$ that is $C^1$-close to the identity. In order to prove this, let us focus on the collar neighborhood $A_C\equiv[0,\rho)\times S^1$ of a connected component~$C$ of $\partial\Sigma$. If we write $(R_t,S_t):=\phi_t(r,s)$, then Equation~\eqref{e:graph_tdF} in the annulus $\mathrm{int}(A_C)$ becomes
\[
R_t\,(S_t-s) = t\,\partial_RF(R_t,s),\qquad
\tfrac 12 (r^2-R_t^2) = t\,\partial_sF(R_t,s),
\]
that is, using \eqref{Fnearbdry},
\[
S_t = s + \underbrace{\big(2 G(R_t,s) + R_t\,\partial_R G(R_t,s) \big) t}_{(*)},\qquad
r = R_t\underbrace{\sqrt{ 1 + 2 t\,\partial_sG(R_t,s) }}_{(**)}.
\]
The term $(*)$ is $C^1$-small and the term $(**)$ is $C^1$-close to 1 as functions of $(R_t,s)$. This shows that the isotopy $(R_t,s)\mto(r,S_t)$ is $C^1$-close to the identity and extends smoothly to the boundary $C$ by $(0,s)\mto(0,s +2 t\,G(0,s))$. Therefore, we obtain a $C^1$-close to the identity smooth extension $\phi_t:\Sigma\to\Sigma$ as well.

Let $X_t$ be the time dependent vector field generating the isotopy $\phi_t$. We claim that~$X_t$ is Hamiltonian with Hamiltonian function
\begin{equation}
\label{H_t}
H_t:\Sigma\to\R,\qquad H_t(z) := F\circ\pi\circ\psi(\phi_t^{-1}(z),z).
\end{equation}
Indeed, consider an arbitrary $v\in\Tan_z\Sigma$ and set
\[
w:=\diff\widetilde\phi_t(z)v, \qquad q_t:=\widetilde\phi_t(z), \qquad y_t:=\tfrac{\diff}{\diff t}\widetilde\phi_t(z).
\]
In Darboux coordinates, we locally see $\Tan^*\Sigma$ as $\Sigma\times\R^2$, and compute
\begin{align*}
\diff\nu\big(X_t(\phi_t(z)),\diff\phi_t(z)v\big)
&=
\psi^*\diff\lambda_{\std} \big( (0,X_t(\phi_t(z))) , (v,\diff\phi_t(z)v) \big)\\
&=
\diff\lambda_{\std}\big( (y_t,\diff F(q_t) + t\,\diff^2F(q_t)y_t) , (w,t\,\diff^2F(q_t)w) \big)\\
&=
\diff F(q_t)w + t\,\diff^2F(q_t)[y_t,w] - t\,\diff^2F(q_t)[w,y_t]\\
&=
\diff F(q_t)w
= \diff (F\circ\widetilde\phi_t)(z)v\\
&= \diff (F\circ\widetilde\phi_t\circ\phi_t^{-1})(\phi_t(z))\diff\phi_t(z)v\\
&=\diff H_t(\phi_t(z))\diff\phi_t(z)v,
\end{align*}
proving our claim.

Note that for every $t\in [0,1]$, the maximum and minimum of $H_t$ on $\Sigma$ coincide with those of $F$. Note also that by \eqref{Fnearbdry} we have
\begin{equation}
\label{sopra}
H_t(z) = F(z) = b_C \qquad \forall z\in C,
\end{equation}
for every connected component $C$ of $\partial \Sigma$. Moreover, the previous considerations on $F$, $b_C$ and $G_C$ imply that if $\phi$ is $C^1$-close to the identity, then $H_t$ is $C^1$-small and on $A_C$ has the form
\[
H_t(r,s) = b_C + r^2 h_C(t,r,s),
\]
where both $|b_C|$ and $\|h_C\|_{C^0}$ are small. Indeed, the above identity and the $C^0$-smallness of $h_C$ follow from \eqref{psi} and \eqref{H_t}, which give us the identity
\[
H_t(r,s) = F(r,\overline{S}_t(r,s)) = b_C + r^2 G_C(r,\overline{S}_t(r,s)),
\]
where $\phi_t^{-1}(r,s)) = ( \overline{R}_t(r,s),\overline{S}_t(r,s))$. Together with the already mentioned $C^1$\nobreakdash-close\-ness of $\phi_t$ to the identity, this proves statements \eqref{t:quasi_autonomousi} and \eqref{t:quasi_autonomousii} in Theorem~\ref{t:quasi_autonomous}.

Let us check that for every $t\in [0,1]$ the function $H_t$ achieves its minimum at some point $z_{\min}\in \Sigma$ which is independent of $t$. If $F$ achieves its minimum on $\partial \Sigma$, this follows from the identity $\min H_t = \min F$ and \eqref{sopra}. So let us assume that $F$ achieves its minimum at an interior point $q_{\min}$. Then $dF(q_{\min})=0$ and, since the inverse image of the zero-section in $T^* \mathrm{int}(\Sigma)$ is the diagonal in $\mathrm{int}(\Sigma) \times \mathrm{int}(\Sigma)$, we~have
\[
\psi(z_{\min},z_{\min}) = (q_{\min},0)
\]
for some $z_{\min} \in \mathrm{int}(\Sigma)$. Since $\psi$ maps the graph of $\phi_t$ to the graph of $t\, dF$, we have
\[
\phi_t(z_{\min}) = z_{\min} \qquad \forall t\in [0,1],
\]
and hence
\[
H_t(z_{\min}) = F \circ \pi \circ \psi(z_{\min},z_{\min}) = F \circ \pi (q_{\min},0) = F(q_{\min}).
\]
This shows that
\[
H_t(z_{\min}) = \min_{\Sigma} H_t \qquad \forall t\in [0,1].
\]
The argument for the maximum is analogous, and we conclude that $H_t$ is quasi-autonomous.
\end{proof}


\begin{thebibliography}{99}
\bibitem{Abbondandolo:2019tl} A. Abbondandolo \& G. Benedetti - “On the local systolic optimality of Zoll contact forms” , 2019
\bibitem{Abbondandolo:2018fb} A. Abbondandolo, B. Bramham, U. L. Hryniewicz \& P. A. S. Salomão - “Sharp systolic inequalities for Reeb flows on the three-sphere” , Invent. Math. 211 (2018) no. 2, p. 687-778
\bibitem{Abbondandolo:2019ta} A. Abbondandolo, B. Bramham, U. L. Hryniewicz \& P. A. S. Salomão - “Contact forms with large systolic ratio in dimension three” , Ann. Scuola Norm. Sup. Pisa Cl. Sci. (5) 19 (2019) no. 4, p. 1561-1582
\bibitem{Albach:2021wi} B. Albach \& H. Geiges - “Surfaces of section for Seifert fibrations” , Arnold Math. J. 7 (2021) no. 4, p. 573-597
\bibitem{Alvarez-Paiva:2014uq} J. C. Álvarez Paiva \& F. Balacheff - “Contact geometry and isosystolic inequalities” , Geom. Funct. Anal. 24 (2014) no. 2, p. 648-669
\bibitem{Banyaga:1978} A. Banyaga - “Sur la structure du groupe des difféomorphismes qui préservent une forme symplectique” , Comment. Math. Helv. 53 (1978) no. 2, p. 174-227
\bibitem{Banyaga:1997} A. Banyaga - The structure of classical diffeomorphism groups , Mathematics and its Appl. , vol. 400 , Kluwer Academic Publishers Group, Dordrecht, 1997
\bibitem{Besse:1978pr} A. L. Besse - Manifolds all of whose geodesics are closed , Ergeb. Math. Grenzgeb. , vol. 93 , Springer-Verlag, Berlin-New York, 1978
\bibitem{Benedetti:2021aa} G. Benedetti \& J. Kang - “A local contact systolic inequality in dimension three” , J. Eur. Math. Soc. (JEMS) 23 (2021) no. 3, p. 721-764
\bibitem{Bottkol:1980} M. Bottkol - “Bifurcation of periodic orbits on manifolds and Hamiltonian systems” , J. Differential Equations 37 (1980) no. 1, p. 12-22
\bibitem{Bialy:1994} M. Bialy \& L. Polterovich - “Geodesics of Hofer’s metric on the group of Hamiltonian diffeomorphisms” , Duke Math. J. 76 (1994) no. 1, p. 273-292
\bibitem{Calabi:1970} E. Calabi - “On the group of automorphisms of a symplectic manifold” , in Problems in analysis , Princeton University Press, Princeton, NJ, 1970, p. 1-26
\bibitem{Cristofaro-Gardiner:2020aa} D. Cristofaro-Gardiner \& M. Mazzucchelli - “The action spectrum characterizes closed contact 3-manifolds all of whose Reeb orbits are closed” , Comment. Math. Helv. 95 (2020) no. 3, p. 461-481
\bibitem{Croke:1994} C. B. Croke \& B. Kleiner - “Conjugacy and rigidity for manifolds with a parallel vector field” , J. Differential Geom. 39 (1994) no. 3, p. 659-680
\bibitem{Epstein:1972} D. B. A. Epstein - “Periodic flows on three-manifolds” , Ann. of Math. (2) 95 (1972), p. 66-82
\bibitem{Geiges:2020aa} H. Geiges - “What does a vector field know about volume?” , J. Fixed Point Theory Appl. 24 (2022) no. 2, article ID 23, 26 pages
\bibitem{Gromoll:1981kl} D. Gromoll \& K. Grove - “On metrics on $S^{2}$ all of whose geodesics are closed” , Invent. Math. 65 (1981/82) no. 1, p. 175-177
\bibitem{Ginzburg:1987lq} V. L. Ginzburg - “New generalizations of Poincaré’s geometric theorem” , Funktsional. Anal. i Prilozhen. 21 (1987) no. 2, p. 16-22
\bibitem{Giroux:2020} E. Giroux - “Ideal Liouville domains, a cool gadget” , J. Symplectic Geom. 18 (2020) no. 3, p. 769-790
\bibitem{Geiges:2018aa} H. Geiges \& C. Lange - “Seifert fibrations of lens spaces” , Abh. Math. Sem. Univ. Hamburg 88 (2018) no. 1, p. 1-22, Correction: Ibid. 91 (2021), no. 1, p. 145–150
\bibitem{Jankins:1983zm} M. Jankins \& W. D. Neumann - Lectures on Seifert manifolds , Brandeis Lect. Notes , vol. 2 , Brandeis University, Waltham, MA, 1983
\bibitem{Katok:1973mw} A. B. Katok - “Ergodic perturbations of degenerate integrable Hamiltonian systems” , Izv. Akad. Nauk SSSR Ser. Mat. 37 (1973), p. 539-576
\bibitem{Kegel:2021vo} M. Kegel \& C. Lange - “A Boothby-Wang theorem for Besse contact manifolds” , Arnold Math. J. 7 (2021) no. 2, p. 225-241
\bibitem{Lange:2020} C. Lange - “On metrics on 2-orbifolds all of whose geodesics are closed” , J. reine angew. Math. 758 (2020), p. 67-94
\bibitem{Lalonde:1995} F. Lalonde \& D. McDuff - “Hofer’s $L^\infty $-geometry: energy and stability of Hamiltonian flows. I, II” , Invent. Math. 122 (1995) no. 1, p. 1-33, 35–69, Erratum: Ibid. 123 (1996), no. 3, p. 613
\bibitem{Lisca:2004oz} P. Lisca \& G. Matić - “Transverse contact structures on Seifert 3-manifolds” , Algebraic Geom. Topol. 4 (2004), p. 1125-1144
\bibitem{Lange:2021} C. Lange \& T. Soethe - “Sharp systolic inequalities for rotationally symmetric $2$-orbifolds” , 2021
\bibitem{Mazzucchelli:2020aa} M. Mazzucchelli \& M. Radeschi - “On the structure of Besse convex contact spheres” , 2020
\bibitem{McDuff:1998uu} D. McDuff \& D. Salamon - Introduction to symplectic topology , Oxford Graduate Texts in Math. , Oxford University Press, Oxford, 2017
\bibitem{Orlik:1972ts} P. Orlik - Seifert manifolds , Lect. Notes in Math. , vol. 291 , Springer-Verlag, Berlin-New York, 1972
\bibitem{Radeschi:2017dz} M. Radeschi \& B. Wilking - “On the Berger conjecture for manifolds all of whose geodesics are closed” , Invent. Math. 210 (2017) no. 3, p. 911-962
\bibitem{Saglam:2018tb} M. Sağlam - “Contact forms with large systolic ratio in arbitrary dimensions” , Ann. Scuola Norm. Sup. Pisa Cl. Sci. (5) 22 (2021) no. 3, p. 1265-1308
\bibitem{Sullivan:1978} D. Sullivan - “A foliation of geodesics is characterized by having no ‘tangent homologies’” , J. Pure Appl. Algebra 13 (1978) no. 1, p. 101-104
\bibitem{Taubes:2007wi} C. H. Taubes - “The Seiberg-Witten equations and the Weinstein conjecture” , Geom. Topol. 11 (2007), p. 2117-2202
\bibitem{Wadsley:1975sp} A. W. Wadsley - “Geodesic foliations by circles” , J. Differential Geom. 10 (1975) no. 4, p. 541-549
\bibitem{Weinstein:1973} A. Weinstein - “Normal modes for nonlinear Hamiltonian systems” , Invent. Math. 20 (1973), p. 47-57
\bibitem{Ziller:1983rw} W. Ziller - “Geometry of the Katok examples” , Ergodic Theory Dynam. Systems 3 (1983) no. 1, p. 135-157
\end{thebibliography}
\end{document}
